% Complete Section 5.2 and Box 5.1, original PDF 64-72 / printed 54-62.
% Source inconsistencies are retained and recorded separately, never repaired here.
% block: mixed-methods-heading page:64 type:section
\subsection{\BT{Mixed methods of discretization}{混合离散方法}}\label{sec:mixed-methods}

% block: mixed-galerkin-heading page:64 type:subsection
\subsubsection{\BT{The mixed Galerkin method}{混合 Galerkin 方法}}\label{subsec:mixed-galerkin}

% block: mixed-galerkin-introduction page:64 type:ordinary-paragraph
\BT{In view of the deficiencies of Galerkin methods based on the primal weak formulation, we now turn to mixed methods, i.e., Galerkin methods based on the \emph{mixed} weak formulation. Here the situation is much better. Indeed, if the Galerkin subspaces are carefully chosen, then a mixed method returns an accurate, stable, convergent approximation, avoiding both of the problems just discussed for the primal formulation. To achieve this, we will impose three conditions on the choice of Galerkin subspaces, and from these we will obtain a variety of results, including stability and convergence.}{鉴于基于原始弱表述的 Galerkin 方法存在上述缺陷, 现在转向混合方法, 即基于\emph{混合}弱表述的 Galerkin 方法. 这里的情况要好得多. 事实上, 如果仔细选取 Galerkin 子空间, 混合方法就能给出准确, 稳定且收敛的近似, 避免刚才讨论的原始表述的两个问题. 为实现这一点, 我们将对 Galerkin 子空间的选取施加三个条件, 并由此得到包括稳定性和收敛性在内的多种结果.}

% block: mixed-discrete-spaces page:64 type:ordinary-paragraph
\BT{Let $V_h^j$ be a finite dimensional subspace of $V^j$. This gives rise to spaces of discrete cocycles and discrete coboundaries:}{设 $V_h^j$ 是 $V^j$ 的有限维子空间. 由此得到离散闭上链空间和离散上边界链空间:}
\[
\mathcal Z_h^j=\{\,v\in V_h^j\mid dv=0\,\}\subset\mathcal Z^j,\qquad\mathcal B_h^{j+1}=\{\,dv\mid v\in V_h^j\,\}\subset\mathcal B^{j+1}.
\]
\BT{Given $V_h^{k-1}\subset V^{k-1}$ and $V_h^k\subset V^k$, the associated space of \emph{discrete harmonic forms} is defined in analogy with \eqref{eq:harmonic-cocycle-orthogonality} as}{给定 $V_h^{k-1}\subset V^{k-1}$ 和 $V_h^k\subset V^k$, 仿照 \eqref{eq:harmonic-cocycle-orthogonality}, 相应的\emph{离散调和形式}空间定义为}
\[
\mathsf H_h^k=\{\,v\in\mathcal Z_h^k\mid v\perp\mathcal B_h^k\,\}.
\]

% block: mixed-discrete-formulation page:64 type:ordinary-paragraph
\BT{Using the spaces $V_h^{k-1}$, $V_h^k$, and $\mathsf H_h^k$, we define the mixed discretization of the Hodge Laplace problem, based on the mixed weak formulation \eqref{eq:hodge-mixed-formulation} in the obvious way. Namely, given $f\in W^k$, we seek $\sigma_h\in V_h^{k-1}$, $u_h\in V_h^k$, and $p_h\in\mathsf H_h^k$ such that}{利用空间 $V_h^{k-1}$, $V_h^k$ 和 $\mathsf H_h^k$, 按照显然的方式, 基于混合弱表述 \eqref{eq:hodge-mixed-formulation} 定义 Hodge Laplace 问题的混合离散化. 即给定 $f\in W^k$, 求 $\sigma_h\in V_h^{k-1}$, $u_h\in V_h^k$ 和 $p_h\in\mathsf H_h^k$, 使得}
\begin{equation}\label{eq:mixed-discrete-formulation}
\begin{aligned}
\langle\sigma_h,\tau\rangle-\langle u_h,d\tau\rangle&=0,&\tau&\in V_h^{k-1},\\
\langle d\sigma_h,v\rangle+\langle du_h,dv\rangle+\langle p_h,v\rangle&=\langle f,v\rangle,&v&\in V_h^k,\\
\langle u_h,q\rangle&=0,&q&\in\mathsf H_h^k.
\end{aligned}
\end{equation}

% block: mixed-nonconformity pages:64,65 type:ordinary-paragraph
\BT{Note that we generally do \emph{not} expect that $\mathsf H_h^j$ is contained in $\mathsf H^j$. (In finite element applications to the Hodge Laplacian, the former will consist of piecewise polynomials and the latter of smooth nonpolynomial functions.) When $\mathsf H_h^j\not\subset\mathsf H^j$, the mixed discretization \eqref{eq:mixed-discrete-formulation} is not a standard Galerkin discretization of \eqref{eq:hodge-mixed-formulation} but rather a generalized Galerkin method. Specifically, the exact solution $u$ of \eqref{eq:hodge-mixed-formulation} need not satisfy the final equation of \eqref{eq:mixed-discrete-formulation}. This is a consistency error, which we will need to account for in our error analysis. In the language of finite elements, \eqref{eq:mixed-discrete-formulation} is, in this case, a \emph{nonconforming} method. Often, however, we have $\mathsf H_h^k=\mathsf H^k=0$. In such cases, the mixed discretization is a conforming Galerkin method.}{注意, 一般\emph{并不}指望 $\mathsf H_h^j$ 包含于 $\mathsf H^j$. (在 Hodge Laplacian 的有限元应用中, 前者由分片多项式构成, 后者则由光滑的非多项式函数构成.) 当 $\mathsf H_h^j\not\subset\mathsf H^j$ 时, 混合离散化 \eqref{eq:mixed-discrete-formulation} 不是 \eqref{eq:hodge-mixed-formulation} 的标准 Galerkin 离散化, 而是广义 Galerkin 方法. 具体地, \eqref{eq:hodge-mixed-formulation} 的精确解 $u$ 不一定满足 \eqref{eq:mixed-discrete-formulation} 的最后一个方程. 这是一种相容性误差, 在误差分析中需要加以考虑. 用有限元的语言说, 此时 \eqref{eq:mixed-discrete-formulation} 是\emph{非协调}方法. 但是, 经常有 $\mathsf H_h^k=\mathsf H^k=0$. 在这种情况下, 混合离散化就是协调 Galerkin 方法.}

% block: mixed-discrete-uniqueness page:65 type:ordinary-paragraph
\BT{The mixed method is well-defined in the sense that \eqref{eq:mixed-discrete-formulation} has a unique solution, no matter how the subspaces $V_h^{k-1}$ and $V_h^k$ are chosen. To prove this, it suffices to show that the only solution for $f=0$ is $\sigma_h=0$, $u_h=0$, $p_h=0$. Choosing $\tau=\sigma_h$, $v=u_h$, and $q=p_h$, and combining the equations, we get that $\sigma_h=0$ and $du_h=0$. Thus $u_h\in\mathcal Z_h$. Moreover, since $\sigma_h=0$, the first equation implies that $u_h\perp\mathcal B_h$. Thus $u_h\in\mathsf H_h$, and so the last equation implies that $u_h$ vanishes. Finally, we take $v=p_h$ in the second equation and conclude that $p_h$ vanishes as well.}{混合方法是良定义的, 即无论如何选取子空间 $V_h^{k-1}$ 和 $V_h^k$, \eqref{eq:mixed-discrete-formulation} 都有唯一解. 为证明这一点, 只需证明在 $f=0$ 时, 唯一解是 $\sigma_h=0$, $u_h=0$, $p_h=0$. 取 $\tau=\sigma_h$, $v=u_h$, $q=p_h$, 并组合这些方程, 得 $\sigma_h=0$ 和 $du_h=0$. 因此 $u_h\in\mathcal Z_h$. 此外, 由于 $\sigma_h=0$, 第一个方程蕴含 $u_h\perp\mathcal B_h$. 因此 $u_h\in\mathsf H_h$, 从而最后一个方程蕴含 $u_h$ 为零. 最后, 在第二个方程中取 $v=p_h$, 可知 $p_h$ 也为零.}

% block: mixed-third-space page:65 type:ordinary-paragraph
\BT{Next we turn to the analysis of this mixed method. Although it is not used in the method itself, for the analysis we require the specification of a third space, $V_h^{k+1}\subset V^{k+1}$.}{接下来分析这一混合方法. 尽管方法本身不使用第三个空间, 分析时仍需要指定 $V_h^{k+1}\subset V^{k+1}$.}

% block: mixed-three-properties-heading page:65 type:subsection
\subsubsection{\BT{Three fundamental properties of the subspaces}{子空间的三个基本性质}}\label{subsec:mixed-three-properties}

% block: mixed-three-properties-intro page:65 type:ordinary-paragraph
\BT{An important realization arising from FEEC is that there are three properties which the subspaces $V_h^j$ should satisfy in order to guarantee that the mixed Galerkin method provides accurate approximation of the true solution. We now specify these properties.}{FEEC 带来的一个重要认识是, 为保证混合 Galerkin 方法准确逼近真解, 子空间 $V_h^j$ 应满足三个性质. 现在具体说明这些性质.}

% block: mixed-approximation-property page:65 type:ordinary-paragraph
\medskip\noindent\textbf{\BT{Approximation property}{逼近性质}}\par
\BT{The first property we need for these spaces is the obviously necessary condition that the subspaces used to approximate $u$ and $\sigma$ provide good approximation. This can be quantified in various ways. For example, we might assume that for $j=k-1$ and $k$, we have a family of spaces $V_h^j$, indexed by $h$ decreasing to $0$, and require that}{这些空间需要满足的第一个性质, 是一个显然必要的条件: 用于逼近 $u$ 和 $\sigma$ 的子空间应提供良好的逼近. 这可以用多种方式量化. 例如, 可以假设对 $j=k-1$ 和 $k$, 有一族空间 $V_h^j$, 由递减至 $0$ 的 $h$ 编号, 并要求}
\begin{equation}\label{eq:galerkin-approximation-property}
\lim_{h\to0}\inf_{v\in V_h^j}\|w-v\|_V=0,\qquad w\in V^j.
\end{equation}
\BT{Or we might require that the best approximation error be $O(h^r)$ for some power $r$ for $w$ belonging to some dense subspace of $V^j$. Rather than pin down a specific requirement, we will make an appropriate approximation assumption when needed.}{也可以要求, 对属于 $V^j$ 的某个稠密子空间的 $w$, 最佳逼近误差对某个幂次 $r$ 为 $O(h^r)$. 我们不预先固定某个具体要求, 而是在需要时作出适当的逼近假设.}

% block: mixed-complex-properties page:65 type:ordinary-paragraph
\BT{Two other properties relate to the complex \eqref{eq:approximation-short-complex}, and not just to the individual spaces.}{另外两个性质涉及复形 \eqref{eq:approximation-short-complex}, 而不仅仅涉及各个单独的空间.}

% block: mixed-subcomplex-property pages:65,66 type:ordinary-paragraph
\medskip\noindent\textbf{\BT{Subcomplex property}{子复形性质}}\par
\BT{We require that $dV_h^{k-1}\subset V_h^k$ and $dV_h^k\subset V_h^{k+1}$ or, in other terminology, that}{要求 $dV_h^{k-1}\subset V_h^k$ 和 $dV_h^k\subset V_h^{k+1}$, 或换一种说法, 要求}
\begin{equation}\label{eq:galerkin-subcomplex}
V_h^{k-1}\xrightarrow{d}V_h^k\xrightarrow{d}V_h^{k+1}
\end{equation}
\BT{is a subcomplex of \eqref{eq:approximation-short-complex}. The discrete complex \eqref{eq:galerkin-subcomplex} then is itself a Hilbert complex. As the base spaces we choose $W_h^j=V_h^j$, but with the $W$ inner product, rather than the $V$ inner product. (Of course all inner products are equivalent on $V_h^j$ since it is finite dimensionsal.) Being finite dimensional, \eqref{eq:galerkin-subcomplex} is a closed Hilbert complex.}{是 \eqref{eq:approximation-short-complex} 的子复形. 此时离散复形 \eqref{eq:galerkin-subcomplex} 本身就是 Hilbert 复形. 基空间取为 $W_h^j=V_h^j$, 但赋予 $W$ 内积, 而不是 $V$ 内积. (当然, 由于 $V_h^j$ 是有限维的, 其上的所有内积都等价.) 由于是有限维的, \eqref{eq:galerkin-subcomplex} 是闭 Hilbert 复形.}

% block: discrete-differential-and-cycles page:66 type:ordinary-paragraph
\BT{As a Hilbert complex in its own right, the discrete complex exhibits all the structures discussed in the previous chapter. The discrete differential}{离散复形本身就是 Hilbert 复形, 因而具有上一章讨论的所有结构. 离散微分}
\[
d_h^j:V_h^j\to V_h^{j+1}
\]
\BT{is just the restriction of $d^j$ (for $j=k-1$ or $k$). The kernel of $d_h^j$ is denoted by $\mathcal Z_h^j$ and the range of $d_h^{j-1}$ by $\mathcal B_h^j$. As for any complex,}{就是 $d^j$ 的限制 (其中 $j=k-1$ 或 $k$). 用 $\mathcal Z_h^j$ 表示 $d_h^j$ 的核, 用 $\mathcal B_h^j$ 表示 $d_h^{j-1}$ 的像. 与任意复形一样, 有}
\[
\mathcal B_h^j\subset\mathcal Z_h^j.
\]
\BT{The orthogonal complement of the former space inside the latter is the space of discrete harmonic forms $\mathsf H_h^k$. Since the discrete differential $d_h^j$ is just a restriction of $d^j$, we have}{前一空间在后一空间中的正交补, 就是离散调和形式空间 $\mathsf H_h^k$. 由于离散微分 $d_h^j$ 只是 $d^j$ 的限制, 有}
\[
\mathcal Z_h^j\subset\mathcal Z^j,\qquad\mathcal B_h^j\subset\mathcal B^j.
\]
\BT{However, it is generally \emph{not true} that $\mathsf H_h^j\subset\mathsf H^j$ since elements of the latter space must be orthogonal to all of $\mathcal B^j$, not just to $\mathcal B_h^j$.}{然而, $\mathsf H_h^j\subset\mathsf H^j$ 一般\emph{不成立}, 因为后一空间的元素必须与整个 $\mathcal B^j$ 正交, 而不只是与 $\mathcal B_h^j$ 正交.}

% block: discrete-adjoint page:66 type:ordinary-paragraph
\BT{The discrete differential $d_h^{j-1}$ has an adjoint operator $d_{jh}^*$. Since $d_h^{j-1}$ is bounded, its adjoint is everywhere defined, so the spaces $V_{jh}^*$ coincide with $W_h^j=V_h^j$ and}{离散微分 $d_h^{j-1}$ 有伴随算子 $d_{jh}^*$. 由于 $d_h^{j-1}$ 有界, 其伴随处处有定义, 所以空间 $V_{jh}^*$ 与 $W_h^j=V_h^j$ 一致, 且}
\[
d_{jh}^*:V_h^j\to V_h^{j-1}.
\]
\BT{Unlike for the discrete exterior derivative itself, the discrete adjoint operator is generally \emph{not} the restriction of the continuous adjoint operator $d_j^*$, and the null spaces $\mathcal Z_{jh}^*$ and ranges $\mathcal B_{jh}^*$ cannot be expected to be contained in their continuous counterparts.}{与离散外微分本身不同, 离散伴随算子一般\emph{不是}连续伴随算子 $d_j^*$ 的限制, 而且不能指望其零空间 $\mathcal Z_{jh}^*$ 和像空间 $\mathcal B_{jh}^*$ 包含于相应的连续空间.}

% block: discrete-hodge-decomposition page:66 type:ordinary-paragraph
\BT{The Hodge decomposition for the discrete complex, which is guaranteed to exist by Theorem~\ref{thm:hodge-decomposition}, is called the \emph{discrete Hodge decomposition}:}{由 Theorem~\ref{thm:hodge-decomposition} 保证存在的离散复形的 Hodge 分解, 称为\emph{离散 Hodge 分解}:}
\begin{equation}\label{eq:discrete-hodge-decomposition}
V_h^k=\mathcal B_h^k\BookOrthSum\mathsf H_h^k\BookOrthSum\mathcal B_{kh}^*.
\end{equation}

% block: bounded-cochain-projection-definition page:66 type:ordinary-paragraph
\medskip\noindent\textbf{\BT{Bounded cochain projections}{有界上链投影}}\par
\BT{Finally, we come to the third and most important assumption, which relates the original complex \eqref{eq:approximation-short-complex} to the subcomplex \eqref{eq:galerkin-subcomplex}. We assume that there exists a \emph{bounded cochain projection} $\pi_h$ from the former to the latter. That it is a \emph{projection} means that for each $j$, $\pi_h^j$ is a linear map from $V^j$ onto $V_h^j$ which restricts to the identity on $V_h^j$. That it is a \emph{cochain} projection means that the following diagram commutes:}{最后讨论第三个也是最重要的假设, 它将原复形 \eqref{eq:approximation-short-complex} 与子复形 \eqref{eq:galerkin-subcomplex} 联系起来. 假设存在从前者到后者的\emph{有界上链投影} $\pi_h$. 所谓\emph{投影}, 是指对每个 $j$, $\pi_h^j$ 都是从 $V^j$ 到 $V_h^j$ 上的线性映射, 且在 $V_h^j$ 上限制为恒等映射. 所谓\emph{上链}投影, 是指下面的图交换:}
\[
\begin{tikzcd}[column sep=large,row sep=large]
V^{k-1}\arrow[r,"d"]\arrow[d,"\pi_h^{k-1}"']&V^k\arrow[r,"d"]\arrow[d,"\pi_h^k"']&V^{k+1}\arrow[d,"\pi_h^{k+1}"]\\
V_h^{k-1}\arrow[r,"d"]&V_h^k\arrow[r,"d"]&V_h^{k+1}
\end{tikzcd}
\]

% block: cochain-boundedness pages:66,67 type:ordinary-paragraph
\BT{Finally, that $\pi_h$ is \emph{bounded} could mean one of two things. It might refer to boundedness in $V$ norm: $\|\pi_hv\|_V\leq c\|v\|_V$ for all $v\in V^j$. Alternatively, we might require $W$-boundedness: $\|\pi_hv\|\leq c\|v\|$. Since $\pi_h$ is a cochain map, $W$-boundedness implies $V$-boundedness:}{最后, $\pi_h$ \emph{有界}可以有两种含义. 它可以指 $V$ 范数下的有界性: 对所有 $v\in V^j$ 有 $\|\pi_hv\|_V\leq c\|v\|_V$. 或者可以要求 $W$-有界性: $\|\pi_hv\|\leq c\|v\|$. 由于 $\pi_h$ 是上链映射, $W$-有界性蕴含 $V$-有界性:}
\[
\begin{aligned}
\|\pi_hv\|_V^2&=\|\pi_hv\|^2+\|d\pi_hv\|^2=\|\pi_hv\|^2+\|\pi_hdv\|^2\\
&\leq c\|v\|^2+\|dv\|^2=c\|v\|_V^2,
\end{aligned}
\]
\BT{where we used $W$-boundedness to obtain the inequality. For now we require only the weaker assumption of $V$-boundedness, although this will change when we consider improved estimates. In either case, if we are considering a family of subcomplexes indexed by $h$, we will want the bounds to be uniform in $h$ in the sense that a single constant $c$ works for all $h$.}{其中用 $W$-有界性得到不等式. 目前只要求较弱的 $V$-有界性假设, 但在考虑改进估计时将改变这一要求. 无论哪种情形, 如果考虑由 $h$ 编号的一族子复形, 都希望界关于 $h$ 一致, 即同一个常数 $c$ 对所有 $h$ 都适用.}

% block: projection-quasioptimality page:67 type:ordinary-paragraph
\BT{A simple, but important, property of a bounded projection is that it provides quasi-optimal approximation, in the sense that, up to a constant multiple, the approximation of an element by its projection is bounded by the best approximation in the space. That is,}{有界投影的一个简单而重要的性质是, 它提供拟最优逼近, 即在相差一个常数倍的意义下, 一个元素的投影逼近误差受该空间中的最佳逼近误差控制. 即}
\begin{equation}\label{eq:projection-quasioptimality}
\begin{aligned}
\|v-\pi_hv\|_V&=\inf_{w\in V_h^j}\|(v-w)-\pi_h(v-w)\|_V\\
&\leq c\inf_{w\in V_h^j}\|v-w\|_V,\qquad v\in V^j,
\end{aligned}
\end{equation}
\BT{with $c=\|I-\pi_h\|_{\mathcal L(V,V)}=\|\pi_h\|_{\mathcal L(V,V)}$. (This last identity goes back to \BookCite{38}. See also \BookCite{69}.) In the rest of this chapter, we derive a series of important further conclusions from these three hypotheses.}{其中 $c=\|I-\pi_h\|_{\mathcal L(V,V)}=\|\pi_h\|_{\mathcal L(V,V)}$. (最后这个恒等式可追溯到 \BookCite{38}. 另见 \BookCite{69}.) 本章余下部分将由这三个假设导出一系列进一步的重要结论.}

% block: mixed-consistency-stability-heading page:67 type:subsection
\subsubsection{\BT{Consistency, stability, and convergence of the mixed method}{混合方法的相容性, 稳定性和收敛性}}\label{subsec:mixed-consistency-stability}

% block: harmonic-gap-heading page:67 type:ordinary-paragraph
\medskip\noindent\textbf{\BT{The gap between $\mathsf H$ and $\mathsf H_h$}{$\mathsf H$ 与 $\mathsf H_h$ 之间的间隙}}\par
\BT{First we show that the cohomology space $\mathcal H_h^k$ for the subcomplex \eqref{eq:galerkin-subcomplex} is isomorphic to that of the original complex \eqref{eq:approximation-short-complex}.}{首先证明子复形 \eqref{eq:galerkin-subcomplex} 的上同调空间 $\mathcal H_h^k$ 同构于原复形 \eqref{eq:approximation-short-complex} 的上同调空间.}

% block: discrete-cohomology-isomorphism page:67 type:theorem
\Needspace{8\baselineskip}
\begin{Theorem}{\BT{isomorphism of cohomology}{上同调的同构}}\label{thm:discrete-cohomology-isomorphism}
\BT{Assume that the complex \eqref{eq:galerkin-subcomplex} is a finite dimensional subcomplex of the \eqref{eq:approximation-short-complex} which admits a $V$-bounded cochain projection $\pi_h$. In addition, assume the approximation property}{假设复形 \eqref{eq:galerkin-subcomplex} 是 \eqref{eq:approximation-short-complex} 的有限维子复形, 且具有 $V$-有界上链投影 $\pi_h$. 此外, 假设逼近性质}
\begin{equation}\label{eq:harmonic-projection-strict-error}
\|q-\pi_hq\|<\|q\|,\qquad0\ne q\in\mathsf H^k.
\end{equation}
\BT{Then $\pi_h$ induces an isomorphism from $\mathcal H^k$ onto $\mathcal H_h^k$.}{则 $\pi_h$ 诱导从 $\mathcal H^k$ 到 $\mathcal H_h^k$ 上的同构.}
\end{Theorem}

% block: discrete-cohomology-isomorphism-proof page:67 type:proof
\begin{Proof}
\BT{As discussed in Section~\ref{sec:chain-maps}, the cochain map $\pi_h$ induces a map $\mathcal H^k\to\mathcal H_h^k$. Since $\pi_h$ is a projection, the induced map on cohomology is a surjection by Proposition~\ref{prop:chain-projection-homology-surjection}. Hence, we need only show that it is an injection. Since $\mathcal H^k$ consists of equivalence classes of elements of $\mathcal Z^k$ modulo $\mathcal B^k$, and $\mathcal H_h^k$ consists of equivalence classes of elements of $\mathcal Z_h^k$ modulo $\mathcal B_h^k$, we need to prove that if $z\in\mathcal Z^k$ and $\pi_hz\in\mathcal B_h^k$, then $z\in\mathcal B^k$. For $z\in\mathcal Z^k$ the Hodge decomposition takes the form $z=q+b$ with $q\in\mathsf H^k$ and $b\in\mathcal B^k$, so our goal is to show that $q$ vanishes. Since $\pi_h$ is a cochain map, $\pi_hb\in\mathcal B_h^k$, and we assumed that $\pi_hz\in\mathcal B_h^k$ as well, so $\pi_hq\in\mathcal B_h^k\subset\mathcal B^k$. But $q$, being harmonic, is orthogonal to $\mathcal B^k$, so $q\perp\pi_hq$, and, by the Pythagorean theorem,}{如 Section~\ref{sec:chain-maps} 所述, 上链映射 $\pi_h$ 诱导映射 $\mathcal H^k\to\mathcal H_h^k$. 由于 $\pi_h$ 是投影, 由 Proposition~\ref{prop:chain-projection-homology-surjection}, 其在上同调上诱导的映射是满射. 因此只需证明它是单射. 由于 $\mathcal H^k$ 由 $\mathcal Z^k$ 中元素模 $\mathcal B^k$ 的等价类构成, $\mathcal H_h^k$ 由 $\mathcal Z_h^k$ 中元素模 $\mathcal B_h^k$ 的等价类构成, 需要证明若 $z\in\mathcal Z^k$ 且 $\pi_hz\in\mathcal B_h^k$, 则 $z\in\mathcal B^k$. 对 $z\in\mathcal Z^k$, Hodge 分解为 $z=q+b$, 其中 $q\in\mathsf H^k$, $b\in\mathcal B^k$, 故目标是证明 $q$ 为零. 因为 $\pi_h$ 是上链映射, 有 $\pi_hb\in\mathcal B_h^k$, 又假设了 $\pi_hz\in\mathcal B_h^k$, 所以 $\pi_hq\in\mathcal B_h^k\subset\mathcal B^k$. 而 $q$ 是调和形式, 与 $\mathcal B^k$ 正交, 所以 $q\perp\pi_hq$, 并且由 Pythagorean 定理,}
\begin{equation}\label{eq:harmonic-projection-pythagoras}
\|q-\pi_hq\|^2=\|q\|^2+\|\pi_hq\|^2\geq\|q\|^2.
\end{equation}
\BT{If $q\ne0$, then \eqref{eq:harmonic-projection-strict-error} and \eqref{eq:harmonic-projection-pythagoras} lead to $\|q\|>\|q\|$, a contradiction, and we conclude that $q=0$ as desired.}{若 $q\ne0$, 则 \eqref{eq:harmonic-projection-strict-error} 和 \eqref{eq:harmonic-projection-pythagoras} 导致 $\|q\|>\|q\|$, 矛盾, 故 $q=0$, 得证.}
\end{Proof}

% block: harmonic-gap-introduction pages:67,68 type:ordinary-paragraph
\BT{Since the homology spaces are isomorphic to the spaces of harmonic forms on both the continuous and the discrete levels, this result tells us that the continuous and discrete harmonic form spaces, $\mathsf H^k$ and $\mathsf H_h^k$, are isomorphic. The next result shows not only that they are isomorphic but also that they are close in the sense that any element of one of the spaces can be well-approximated by an element of the other. This will enable us to control the consistency error coming from the nonconformity of the mixed method when the two spaces are not equal.}{由于在连续和离散层次上, 同调空间都与调和形式空间同构, 此结果说明连续和离散调和形式空间 $\mathsf H^k$ 与 $\mathsf H_h^k$ 同构. 下一个结果不仅表明它们同构, 还说明它们相近, 即其中任一空间的任意元素都能被另一空间的元素良好逼近. 当两个空间不相等时, 这将使我们能够控制混合方法的非协调性所带来的相容性误差.}

% block: harmonic-gap-theorem page:68 type:theorem
\Needspace{8\baselineskip}
\begin{Theorem}{\BT{gap between harmonic forms}{调和形式之间的间隙}}\label{thm:harmonic-gap}
\BT{Assume that the complex \eqref{eq:galerkin-subcomplex} is a finite dimensional subcomplex of \eqref{eq:approximation-short-complex} which admits a $V$-bounded cochain projection $\pi_h$. Then}{假设复形 \eqref{eq:galerkin-subcomplex} 是 \eqref{eq:approximation-short-complex} 的有限维子复形, 且具有 $V$-有界上链投影 $\pi_h$. 则}
\begin{equation}\label{eq:harmonic-continuous-to-discrete-gap}
\|(I-P_{\mathsf H_h})q\|_V\leq\|(I-\pi_h^k)q\|_V,\qquad q\in\mathsf H^k,
\end{equation}
\begin{equation}\label{eq:harmonic-discrete-to-continuous-gap}
\|(I-P_{\mathsf H})q\|_V\leq\|(I-\pi_h^k)P_{\mathsf H}q\|_V,\qquad q\in\mathsf H_h^k.
\end{equation}
\end{Theorem}

% block: harmonic-gap-proof page:68 type:proof
\begin{Proof}
\BT{For $q\in\mathsf H$, $P_{\mathsf H_h}q=P_{\mathcal Z_h}q$, since $\mathcal Z_h=\mathsf H_h\BookOrthSum\mathcal B_h$ and $q\perp\mathcal B\supset\mathcal B_h$. This means that $P_{\mathsf H_h}q$ is the element of $\mathcal Z_h$ closest to $q$ as measured in the $V$ norm. But $\pi_hq$ also belongs to $\mathcal Z_h$, since $\pi_h$ is a cochain map. Therefore, \eqref{eq:harmonic-continuous-to-discrete-gap} holds.}{对 $q\in\mathsf H$, 有 $P_{\mathsf H_h}q=P_{\mathcal Z_h}q$, 因为 $\mathcal Z_h=\mathsf H_h\BookOrthSum\mathcal B_h$ 且 $q\perp\mathcal B\supset\mathcal B_h$. 这意味着, 按 $V$ 范数衡量, $P_{\mathsf H_h}q$ 是 $\mathcal Z_h$ 中离 $q$ 最近的元素. 而 $\pi_hq$ 也属于 $\mathcal Z_h$, 因为 $\pi_h$ 是上链映射. 因此 \eqref{eq:harmonic-continuous-to-discrete-gap} 成立.}

\BT{If $q\in\mathsf H_h\subset\mathcal Z_h\subset\mathcal Z$, the Hodge decomposition gives us $q-P_{\mathsf H}q\in\mathcal B$, so $\pi_h(q-P_{\mathsf H}q)\in\mathcal B_h$, and so is orthogonal to both the discrete harmonic form $q$ and the harmonic form $P_{\mathsf H}q$. Therefore}{若 $q\in\mathsf H_h\subset\mathcal Z_h\subset\mathcal Z$, 由 Hodge 分解可得 $q-P_{\mathsf H}q\in\mathcal B$, 所以 $\pi_h(q-P_{\mathsf H}q)\in\mathcal B_h$, 因而它既与离散调和形式 $q$ 正交, 也与调和形式 $P_{\mathsf H}q$ 正交. 因此}
\[
\begin{aligned}
\|q-P_{\mathsf H}q\|_V&\leq\|q-P_{\mathsf H}q-\pi_h(q-P_{\mathsf H}q)\|_V\\
&=\|(I-\pi_h)P_{\mathsf H}q\|_V.
\end{aligned}
\]
\end{Proof}

% block: subspace-gap-definition page:68 type:ordinary-paragraph
\BT{The \emph{gap} between two closed subspaces $E$ and $F$ of a Hilbert space $V$ is defined [\BookNumber{56}, Chapter IV, Section 2.1]}{Hilbert 空间 $V$ 的两个闭子空间 $E$ 和 $F$ 之间的\emph{间隙}定义为 [\BookNumber{56}, Chapter IV, Section 2.1]}
\begin{equation}\label{eq:subspace-gap-definition}
\operatorname{gap}(E,F)=\max\left(\sup_{\substack{u\in E\\\|u\|=1}}\inf_{v\in F}\|u-v\|_V,\ \sup_{\substack{v\in F\\\|v\|=1}}\inf_{u\in E}\|u-v\|_V\right).
\end{equation}
\BT{An equivalent expression is $\operatorname{gap}(E,F)=\|P_E-P_F\|$, where $P_E,P_F:V\to V$ are the orthogonal projections onto $E$ and $F$ [\BookNumber{56}, footnote 1, p. 198]. From Theorem~\ref{thm:harmonic-gap}, we immediately get an estimate of the gap between the spaces of harmonic and discrete harmonic forms:}{一个等价表达式是 $\operatorname{gap}(E,F)=\|P_E-P_F\|$, 其中 $P_E,P_F:V\to V$ 分别为到 $E$ 和 $F$ 上的正交投影 [\BookNumber{56}, footnote 1, p. 198]. 由 Theorem~\ref{thm:harmonic-gap}, 立即得到调和形式空间与离散调和形式空间之间间隙的估计:}
\[
\operatorname{gap}(\mathsf H,\mathsf H_h)\leq\sup_{\substack{q\in\mathsf H\\\|q\|=1}}\|(I-\pi_h)q\|_V.
\]

% block: discrete-stability-introduction page:68 type:ordinary-paragraph
\medskip\noindent\textbf{\BT{The discrete Poincaré inequality and stability}{离散 Poincaré 不等式与稳定性}}\par
\BT{The \emph{stability} of the mixed method \eqref{eq:mixed-discrete-formulation} is measured by the norm of the discrete solution operator $f\mapsto(\sigma_h,u_h,p_h)$. Since \eqref{eq:mixed-discrete-formulation} is just the mixed weak formulation of the abstract Hodge Laplacian problem associated to the discrete complex \eqref{eq:galerkin-subcomplex}, Theorems~\ref{thm:hodge-formulation-equivalence} and~\ref{thm:hodge-well-posedness}, which apply to any closed Hilbert complex, give the estimate}{混合方法 \eqref{eq:mixed-discrete-formulation} 的\emph{稳定性}, 用离散解算子 $f\mapsto(\sigma_h,u_h,p_h)$ 的范数衡量. 由于 \eqref{eq:mixed-discrete-formulation} 正是与离散复形 \eqref{eq:galerkin-subcomplex} 相伴的抽象 Hodge Laplacian 问题的混合弱表述, 对任意闭 Hilbert 复形都适用的 Theorem~\ref{thm:hodge-formulation-equivalence} 和 Theorem~\ref{thm:hodge-well-posedness} 给出估计}
\begin{equation}\label{eq:discrete-mixed-stability}
\|\sigma_h\|_V+\|u_h\|_V+\|p_h\|\leq c_h\|f\|,
\end{equation}
\BT{with the stability constant $c_h$ depending only on the constant in the Poincaré inequality for the discrete complex. Thus the key to stability is a bound on the discrete Poincaré constant. This is the content of the next result, which gives such a bound in terms of the Poincaré constant for the original complex and the bound on the cochain projection.}{其中稳定性常数 $c_h$ 仅依赖于离散复形的 Poincaré 不等式中的常数. 因此, 稳定性的关键是离散 Poincaré 常数的界. 下一个结果正给出这一界, 用原复形的 Poincaré 常数和上链投影的界来表示.}

% block: discrete-poincare-theorem page:69 type:theorem
\Needspace{8\baselineskip}
\begin{Theorem}{\BT{discrete Poincaré inequality}{离散 Poincaré 不等式}}\label{thm:discrete-poincare}
\BT{Assume that the complex \eqref{eq:galerkin-subcomplex} is a finite dimensional subcomplex of \eqref{eq:approximation-short-complex} which admits a $V$-bounded cochain projection $\pi_h$. Then}{假设复形 \eqref{eq:galerkin-subcomplex} 是 \eqref{eq:approximation-short-complex} 的有限维子复形, 且具有 $V$-有界上链投影 $\pi_h$. 则}
\[
\|v\|_V\leq c_P\|\pi_h^k\|\|dv\|_V,\qquad v\in\mathcal Z_h^{k\perp}\cap V_h^k,
\]
\BT{where $c_P$ is the constant appearing in the Poincaré inequality \eqref{eq:hilbert-poincare-inequality} and $\|\pi_h^k\|$ denotes the $V^k$ operator norm of $\pi_h^k$.}{其中 $c_P$ 是 Poincaré 不等式 \eqref{eq:hilbert-poincare-inequality} 中的常数, $\|\pi_h^k\|$ 表示 $\pi_h^k$ 的 $V^k$ 算子范数.}
\end{Theorem}

% block: discrete-poincare-proof page:69 type:proof
\begin{Proof}
\BT{Given $v\in\mathcal Z_h^{k\perp}\cap V_h^k$, define $z\in\mathcal Z^{k\perp}\cap V^k$ by $dz=dv$. By Theorem~\ref{thm:hilbert-complex-poincare}, $\|z\|_V\leq c_P\|dv\|$. Now, $v-\pi_hz\in V_h^k$ and $d\pi_hz=\pi_hdz=\pi_hdv=dv$ with the three equalities coming from, first, the fact that $\pi_h$ is a cochain map, second, the definition of $z$, and, third, the subcomplex property $dV_h^k\subset V_h^{k+1}$ and the fact that $\pi_h^{k+1}$ is a projection. Thus $v-\pi_hz\in\mathcal Z_h^k$ and so is orthogonal to $v$. Therefore}{给定 $v\in\mathcal Z_h^{k\perp}\cap V_h^k$, 由 $dz=dv$ 定义 $z\in\mathcal Z^{k\perp}\cap V^k$. 由 Theorem~\ref{thm:hilbert-complex-poincare}, 有 $\|z\|_V\leq c_P\|dv\|$. 现在 $v-\pi_hz\in V_h^k$, 且 $d\pi_hz=\pi_hdz=\pi_hdv=dv$; 这三个等式依次来自 $\pi_h$ 是上链映射, $z$ 的定义, 以及子复形性质 $dV_h^k\subset V_h^{k+1}$ 和 $\pi_h^{k+1}$ 为投影这一事实. 因此 $v-\pi_hz\in\mathcal Z_h^k$, 从而与 $v$ 正交. 故}
\[
\begin{aligned}
\|v\|_V^2&=\langle v,\pi_h^kz\rangle_V+\langle v,v-\pi_h^kz\rangle_V\\
&=\langle v,\pi_h^kz\rangle_V\leq\|v\|_V\|\pi_h^k\|\|z\|_V,
\end{aligned}
\]
\BT{so}{因此}
\[
\|v\|_V\leq\|\pi_h^k\|\|z\|_V\leq\|\pi_h^k\|c_P\|dv\|_V.
\]
\end{Proof}

% block: discrete-inf-sup-introduction page:69 type:ordinary-paragraph
\BT{As an immediate consequence of this result and Theorem~\ref{thm:hodge-mixed-inf-sup}, we obtain the discrete inf-sup condition and the stability estimate.}{作为此结果和 Theorem~\ref{thm:hodge-mixed-inf-sup} 的直接推论, 得到离散 inf-sup 条件和稳定性估计.}

% block: discrete-inf-sup-theorem page:69 type:theorem
\Needspace{8\baselineskip}
\begin{Theorem}\label{thm:discrete-mixed-inf-sup}
\BT{Assume that the complex \eqref{eq:galerkin-subcomplex} is a finite dimensional subcomplex of the \eqref{eq:approximation-short-complex} which admits a $V$-bounded cochain projection $\pi_h$. Then the inf-sup condition for the mixed bilinear form $B$, given in \eqref{eq:hodge-mixed-bilinear-form}, over the subspace $V_h^{k-1}\times V_h^k\times\mathsf H_h^k$ holds, and the inf-sup constant $\gamma_h$ can be bounded below by a positive constant depending only on $c_P$ and $\|\pi_h^k\|$. Consequently the stability estimate \eqref{eq:discrete-mixed-stability} holds with a constant $c_h$ depending only on the same quantities.}{假设复形 \eqref{eq:galerkin-subcomplex} 是 \eqref{eq:approximation-short-complex} 的有限维子复形, 且具有 $V$-有界上链投影 $\pi_h$. 则 \eqref{eq:hodge-mixed-bilinear-form} 给出的混合双线性型 $B$ 在子空间 $V_h^{k-1}\times V_h^k\times\mathsf H_h^k$ 上满足 inf-sup 条件, 且 inf-sup 常数 $\gamma_h$ 具有仅依赖于 $c_P$ 和 $\|\pi_h^k\|$ 的正常数下界. 从而稳定性估计 \eqref{eq:discrete-mixed-stability} 成立, 其中常数 $c_h$ 也仅依赖于这些量.}
\end{Theorem}

% block: mixed-basic-error-introduction page:69 type:ordinary-paragraph
\medskip\noindent\textbf{\BT{Error estimates for the mixed method}{混合方法的误差估计}}\par
\BT{In the numerical analysis of differential equations, error estimates are obtained from the stability and consistency of the discretization scheme. In the present case, stability comes from Theorem~\ref{thm:discrete-mixed-inf-sup} and consistency comes from the assumed approximation properties of the space $V_h^j$ and from Theorem~\ref{thm:harmonic-gap} on the gap between the spaces of continuous and discrete harmonic forms (the latter required in the case $\mathsf H_h^k$ is not contained in $\mathsf H^k$, so the mixed method is nonconforming). The following theorem gives the resulting error estimates. It was proven in the concrete case of the de Rham complex in [\BookNumber{11}, Theorem 7.4] and then in the case of general Hilbert complexes (essentially with the same proof) in [\BookNumber{13}, Theorem 3.9]. We repeat this proof here.}{在微分方程的数值分析中, 误差估计由离散格式的稳定性和相容性得到. 在当前情形中, 稳定性来自 Theorem~\ref{thm:discrete-mixed-inf-sup}, 相容性来自对空间 $V_h^j$ 所作的逼近性质假设, 以及关于连续和离散调和形式空间之间间隙的 Theorem~\ref{thm:harmonic-gap} (当 $\mathsf H_h^k$ 不包含于 $\mathsf H^k$, 因而混合方法非协调时, 需要后一个结果). 下述定理给出了由此得到的误差估计. 文献 [\BookNumber{11}, Theorem 7.4] 在 de Rham 复形的具体情形中证明了它, 随后文献 [\BookNumber{13}, Theorem 3.9] 在一般 Hilbert 复形的情形中给出了本质上相同的证明. 这里重复该证明.}

% block: mixed-basic-error-theorem pages:69,70 type:theorem
\Needspace{8\baselineskip}
\begin{Theorem}{\BT{basic error estimate for the mixed method}{混合方法的基本误差估计}}\label{thm:mixed-basic-error}
\BT{Assume that the complex \eqref{eq:galerkin-subcomplex} is a finite dimensional subcomplex of \eqref{eq:approximation-short-complex} which admits a $V$-bounded cochain projection $\pi_h$. Let $(\sigma,u,p)\in V^{k-1}\times V^k\times\mathsf H^k$ solve the Hodge Laplace problem \eqref{eq:hodge-mixed-formulation} and let $(\sigma_h,u_h,p_h)\in V_h^{k-1}\times V_h^k\times\mathsf H_h^k$ solve the discrete problem \eqref{eq:mixed-discrete-formulation}. Then}{假设复形 \eqref{eq:galerkin-subcomplex} 是 \eqref{eq:approximation-short-complex} 的有限维子复形, 且具有 $V$-有界上链投影 $\pi_h$. 设 $(\sigma,u,p)\in V^{k-1}\times V^k\times\mathsf H^k$ 求解 Hodge Laplace 问题 \eqref{eq:hodge-mixed-formulation}, $(\sigma_h,u_h,p_h)\in V_h^{k-1}\times V_h^k\times\mathsf H_h^k$ 求解离散问题 \eqref{eq:mixed-discrete-formulation}. 则}
\[
\begin{aligned}
&\|\sigma-\sigma_h\|_V+\|u-u_h\|_V+\|p-p_h\|\\
&\quad\leq C\Big(\inf_{\tau\in V_h^{k-1}}\|\sigma-\tau\|_V+\inf_{v\in V_h^k}\|u-v\|_V\\
&\hspace{45mm}+\inf_{q\in V_h^k}\|p-q\|_V+\mu\inf_{v\in V_h^k}\|P_{\mathcal B}u-v\|_V\Big),
\end{aligned}
\]
\BT{where}{其中}
\[
\mu=\mu_h^k=\sup_{\substack{r\in\mathsf H^k\\\|r\|=1}}\|(I-\pi_h^k)r\|
\]
\BT{and the constant $C$ depends only on $c_P$ and $\|\pi_h^k\|$.}{且常数 $C$ 仅依赖于 $c_P$ 和 $\|\pi_h^k\|$.}
\end{Theorem}

% block: mixed-basic-error-proof page:70 type:proof
\begin{Proof}
\BT{With $B$ defined as in \eqref{eq:hodge-mixed-bilinear-form}, we first observe that $(\sigma,u,p)$ satisfies}{将 $B$ 定义为 \eqref{eq:hodge-mixed-bilinear-form} 中的双线性型, 首先注意到 $(\sigma,u,p)$ 满足}
\[
\begin{gathered}
B(\sigma,u,p;\tau_h,v_h,q_h)=\langle f,v_h\rangle-\langle u,q_h\rangle,\\
(\tau_h,v_h,q_h)\in V_h^{k-1}\times V_h^k\times\mathsf H_h^k.
\end{gathered}
\]
\BT{Let $\tau$, $v$, and $q$ be the $V$-orthogonal projections of $\sigma$, $u$, and $p$ into $V_h^{k-1}$, $V_h^k$, and $\mathsf H_h^k$, respectively. Then, for any $(\tau_h,v_h,q_h)\in V_h^{k-1}\times V_h^k\times\mathsf H_h^k$, we have}{设 $\tau$, $v$ 和 $q$ 分别为 $\sigma$, $u$ 和 $p$ 到 $V_h^{k-1}$, $V_h^k$ 和 $\mathsf H_h^k$ 中的 $V$-正交投影. 则对任意 $(\tau_h,v_h,q_h)\in V_h^{k-1}\times V_h^k\times\mathsf H_h^k$, 有}
\[
\begin{aligned}
&B(\sigma_h-\tau,u_h-v,p_h-q;\tau_h,v_h,q_h)\\
&\quad=B(\sigma-\tau,u-v,p-q;\tau_h,v_h,q_h)+\langle u,q_h\rangle\\
&\quad=B(\sigma-\tau,u-v,p-q;\tau_h,v_h,q_h)+\langle P_{\mathsf H_h}u,q_h\rangle\\
&\quad\leq C(\|\sigma-\tau\|_V+\|u-v\|_V+\|p-q\|+\|P_{\mathsf H_h}u\|)\\
&\hspace{55mm}(\|\tau_h\|_V+\|v_h\|_V+\|q_h\|).
\end{aligned}
\]
\BT{Theorem~\ref{thm:discrete-mixed-inf-sup} then gives}{于是 Theorem~\ref{thm:discrete-mixed-inf-sup} 给出}
\begin{equation}\label{eq:mixed-error-discrete-stability}
\begin{aligned}
&\|\sigma_h-\tau\|_V+\|u_h-v\|_V+\|p_h-q\|\\
&\qquad\leq C(\|\sigma-\tau\|_V+\|u-v\|_V+\|p-q\|+\|P_{\mathsf H_h}u\|).
\end{aligned}
\end{equation}
\BT{Using \eqref{eq:harmonic-continuous-to-discrete-gap} and the boundedness of the projection $\pi_h$ we have}{利用 \eqref{eq:harmonic-continuous-to-discrete-gap} 和投影 $\pi_h$ 的有界性, 有}
\begin{equation}\label{eq:mixed-error-harmonic-best-approximation}
\|p-q\|\leq\|(I-\pi_h)p\|\leq C\inf_{q\in V_h^k}\|p-q\|_V.
\end{equation}
\BT{Next we show that}{接下来证明}
\[
\|P_{\mathsf H_h}u\|\leq\mu\|(I-\pi_h)u_{\mathcal B}\|_V.
\]
\BT{Now $u\perp\mathsf H^k$, so $u=u_{\mathcal B}+u_\perp$, with $u_{\mathcal B}\in\mathcal B^k$ and $u_\perp\in\mathcal Z^{k\perp}$. Since $\mathsf H_h^k\subset\mathcal Z^k$, $P_{\mathsf H_h}u_\perp=0$, and since $\pi_hu_{\mathcal B}\in\mathcal B_h^k$, $P_{\mathsf H_h}\pi_hu_{\mathcal B}=0$. Let $q=P_{\mathsf H_h}u/\|P_{\mathsf H_h}u\|\in\mathsf H_h^k$. By Theorem~\ref{thm:harmonic-gap}, there exists $r\in\mathsf H^k$ (and so $r\perp\mathcal B^k$) with $\|r\|\leq1$ and}{由于 $u\perp\mathsf H^k$, 有 $u=u_{\mathcal B}+u_\perp$, 其中 $u_{\mathcal B}\in\mathcal B^k$, $u_\perp\in\mathcal Z^{k\perp}$. 由于 $\mathsf H_h^k\subset\mathcal Z^k$, 有 $P_{\mathsf H_h}u_\perp=0$; 又由于 $\pi_hu_{\mathcal B}\in\mathcal B_h^k$, 有 $P_{\mathsf H_h}\pi_hu_{\mathcal B}=0$. 令 $q=P_{\mathsf H_h}u/\|P_{\mathsf H_h}u\|\in\mathsf H_h^k$. 由 Theorem~\ref{thm:harmonic-gap}, 存在 $r\in\mathsf H^k$ (因而 $r\perp\mathcal B^k$), 满足 $\|r\|\leq1$ 以及}
\[
\|q-r\|\leq\|(I-\pi_h)r\|\leq\sup_{\substack{r\in\mathsf H^k\\\|r\|=1}}\|(I-\pi_h)r\|.
\]
\BT{Therefore}{因此}
\begin{equation}\label{eq:mixed-error-consistency-bound}
\begin{aligned}
\|P_{\mathsf H_h}u\|&=\langle u_{\mathcal B}-\pi_hu_{\mathcal B},q-r\rangle\\
&\leq\|(I-\pi_h)u_{\mathcal B}\|\sup_{\substack{r\in\mathsf H^k\\\|r\|=1}}\|(I-\pi_h)r\|\\
&\leq c\mu\inf_{v\in V_h^k}\|P_{\mathcal B}u-v\|_V),
\end{aligned}
\end{equation}
\BT{since $\pi_h$ is a bounded projection. The theorem follows from \eqref{eq:mixed-error-discrete-stability}--\eqref{eq:mixed-error-consistency-bound} and the triangle inequality.}{因为 $\pi_h$ 是有界投影. 由 \eqref{eq:mixed-error-discrete-stability}--\eqref{eq:mixed-error-consistency-bound} 及三角不等式即得本定理.}
\end{Proof}

% block: mixed-basic-error-interpretation page:71 type:ordinary-paragraph
\BT{Theorem~\ref{thm:mixed-basic-error} is typically applied when we have a family of Galerkin subspaces $V_h^j$ indexed by $h$ (usually a mesh size parameter), with uniformly bounded cochain projections. Then the rate of the convergence depends only on approximation properties of the spaces $V_h^{k-1}$ and $V_h^k$. It is important to note that the third term on the right-hand side of the error estimate is the approximation for $p$ by the full space $V_h^k$, not just by $\mathsf H_h^k$. The final term on the right-hand side is typically much smaller than the other terms. If there are no harmonic forms, it is not present at all. Otherwise, it is the product of two terms, both of which are determined by the approximation properties of $V_h^k$.}{Theorem~\ref{thm:mixed-basic-error} 通常用于这样一种情形: 有一族由 $h$ (通常为网格尺寸参数) 编号的 Galerkin 子空间 $V_h^j$, 且上链投影一致有界. 此时收敛阶只依赖于空间 $V_h^{k-1}$ 和 $V_h^k$ 的逼近性质. 需要注意, 误差估计右端第三项是用整个空间 $V_h^k$ 逼近 $p$, 而不只是用 $\mathsf H_h^k$. 右端最后一项通常远小于其他项. 如果不存在调和形式, 这一项就完全不出现. 否则, 它是两项的乘积, 两项都由 $V_h^k$ 的逼近性质决定.}

% block: mixed-sharper-error-analogy page:71 type:ordinary-paragraph
\BT{The left-hand side of the error estimate couples together the $W$ norm of various solution quantities: $\sigma$, $d\sigma$, $u$, $du$, and $p$. It is possible to get sharper estimates by considering them individually. An analogous situation (actually a special case) is the error estimate for the standard piecewise linear finite element method for Poisson's equation, which takes the form}{误差估计左端耦合了各种解变量的 $W$ 范数: $\sigma$, $d\sigma$, $u$, $du$ 和 $p$. 分别考虑这些量, 可能得到更精细的估计. 一个类似情形 (实际上是特殊情形) 是 Poisson 方程的标准分片线性有限元方法的误差估计, 其形式为}
\[
\|u-u_h\|_{H^1}\leq C\inf_{v\in V_h}\|u-v\|_{H^1}\leq Ch\|u\|_{H^2}.
\]
\BT{The left-hand side combines the $L^2$ error of $u$ with that of $\operatorname{grad}u$ and so gives a suboptimal rate for the $L^2$ error in $u$. If we apply the Aubin–Nitsche duality argument, however, we can prove the sharper estimate $\|u-u_h\|\leq Ch^2\|u\|_{H^2}$.}{左端将 $u$ 的 $L^2$ 误差与 $\operatorname{grad}u$ 的误差结合在一起, 因而对 $u$ 的 $L^2$ 误差给出次优阶. 但是, 如果应用 Aubin–Nitsche 对偶论证, 就能证明更精细的估计 $\|u-u_h\|\leq Ch^2\|u\|_{H^2}$.}

% block: improved-dsigma-example page:71 type:ordinary-paragraph
\BT{As a simple example of such an improved error estimates for the mixed Hodge Laplacian, consider the error in $d\sigma$ measured in the norm of $W^k$. From the relation between the mixed formulation and the Hodge decomposition, we have}{作为混合 Hodge Laplacian 的此类改进误差估计的一个简单例子, 考虑用 $W^k$ 范数衡量的 $d\sigma$ 的误差. 由混合表述与 Hodge 分解之间的关系, 有}
\[
d\sigma_h=P_{\mathcal B_h}f=P_{\mathcal B_h}P_{\mathcal B}f=P_{\mathcal B_h}d\sigma.
\]
\BT{Since $P_{\mathcal B_h}d\sigma$ is the best approximation of $d\sigma$ from $\mathcal B_h$ and $\pi_hd\sigma$ is another element of $\mathcal B_h$, we have}{由于 $P_{\mathcal B_h}d\sigma$ 是 $d\sigma$ 在 $\mathcal B_h$ 中的最佳逼近, 且 $\pi_hd\sigma$ 是 $\mathcal B_h$ 的另一个元素, 有}
\[
\|d\sigma-d\sigma_h\|=\|d\sigma-P_{\mathcal B_h}d\sigma\|\leq\|d\sigma-\pi_hd\sigma\|.
\]
\BT{Assuming that the $\pi_h^k$ is $W$-bounded, we obtain}{假设 $\pi_h^k$ 是 $W$-有界的, 得到}
\[
\|d\sigma-d\sigma_h\|\leq\|\pi_h^k\|_{\mathcal L(W,W)}\inf_{v\in V_h^k}\|d\sigma-v\|,
\]
\BT{just as in \eqref{eq:projection-quasioptimality}. Thus, approximation of $d\sigma$ by $d\sigma_h$ is, up to a constant factor, as good as the best approximation afforded by the subspace $V_h^k$ where $d\sigma_h$ lives.}{与 \eqref{eq:projection-quasioptimality} 一样. 因此, 在相差一个常数因子的意义下, 用 $d\sigma_h$ 逼近 $d\sigma$, 与 $d\sigma_h$ 所在子空间 $V_h^k$ 提供的最佳逼近同样好.}

% block: improved-error-introduction page:71 type:ordinary-paragraph
\BT{Theorem~\ref{thm:mixed-improved-errors} below, which is taken from \BookCite{13}, gives the full set of such improved error estimates for the mixed method for the Hodge Laplacian. The theorem requires not just that the Hilbert complex on the continuous level be closed but rather that it satisfy the compactness property (see Definition~\ref{def:hilbert-complex-compactness-property}), which is a stronger assumption. Moreover, it requires that the discrete subcomplex admit not just $V$-bounded cochain projections but $W$-bounded. The theorem is stated in terms of the best approximation error in the $W^j$ norm by elements of $V_h^j$ for various quantities $w\in W^j$, for which we introduce the notation}{下面取自 \BookCite{13} 的 Theorem~\ref{thm:mixed-improved-errors}, 给出了 Hodge Laplacian 混合方法的此类改进误差估计的完整集合. 该定理不仅要求连续层次上的 Hilbert 复形是闭的, 而且要求它满足紧性性质 (见 Definition~\ref{def:hilbert-complex-compactness-property}), 这是更强的假设. 此外, 它要求离散子复形具有的上链投影不只是 $V$-有界, 而且是 $W$-有界的. 该定理用各种量 $w\in W^j$ 被 $V_h^j$ 中元素在 $W^j$ 范数下的最佳逼近误差来表述, 为此引入记号}
\[
E(w)=E_h^j(w)=\inf_{v\in V_h^j}\|w-v\|,\qquad w\in W^j.
\]

% block: improved-error-parameters pages:71,72 type:ordinary-paragraph
\BT{Finally, the statement of the theorem uses the notation}{最后, 定理的陈述还使用记号}
\[
\begin{aligned}
\mu=\mu_h^k&=\|(I-\pi_h)P_{\mathsf H}\|_{\mathcal L(W,W)},\\
\delta=\delta_h^k&=\|(I-\pi_h)K\|_{\mathcal L(W,W)},\\
\eta=\eta_h^k&=\max_{j=k-1,k}\big[\|(I-\pi_h^j)d^*K\|_{\mathcal L(W,W)},\ \|(I-\pi_h^{j+1})dK\|_{\mathcal L(W,W)}\big],
\end{aligned}
\]
\BT{of which the first already appeared in Theorem~\ref{thm:mixed-basic-error}. The operator $K$ here is the solution operator for the Hodge Laplacian defined just after Remark~\ref{rem:hodge-general-mixed-load}. Due to the compactness property, $K$, $dK$, and $d^*K$ are compact operators on $W$. Together with an approximation property, which ensures that $\pi_h$ converges pointwise to the identity on $W$ as $h\to0$, this implies that $\eta,\delta,\mu\to0$ with $h$.}{其中第一个已经在 Theorem~\ref{thm:mixed-basic-error} 中出现. 这里的算子 $K$ 是紧接 Remark~\ref{rem:hodge-general-mixed-load} 之后定义的 Hodge Laplacian 的解算子. 由紧性性质, $K$, $dK$ 和 $d^*K$ 都是 $W$ 上的紧算子. 再结合保证 $\pi_h$ 在 $h\to0$ 时逐点收敛于 $W$ 上恒等映射的逼近性质, 可得 $\eta,\delta,\mu\to0$, 随 $h$ 趋于零.}

% block: mixed-improved-error-theorem page:72 type:theorem
\Needspace{8\baselineskip}
\begin{Theorem}{\BT{improved error estimates for the mixed method}{混合方法的改进误差估计}}\label{thm:mixed-improved-errors}
\BT{Assume that the Hilbert complex \eqref{eq:approximation-short-complex} satisfies the compactness property and that \eqref{eq:galerkin-subcomplex} is a finite dimensional subcomplex which admits a $W$-bounded cochain projection $\pi_h$. Let $(\sigma,u,p)\in V^{k-1}\times V^k\times\mathsf H^k$ solve the Hodge Laplace problem \eqref{eq:hodge-mixed-formulation} and let $(\sigma_h,u_h,p_h)\in V_h^{k-1}\times V_h^k\times\mathsf H_h^k$ solve the discrete problem \eqref{eq:mixed-discrete-formulation}. Then}{假设 Hilbert 复形 \eqref{eq:approximation-short-complex} 满足紧性性质, 且 \eqref{eq:galerkin-subcomplex} 是具有 $W$-有界上链投影 $\pi_h$ 的有限维子复形. 设 $(\sigma,u,p)\in V^{k-1}\times V^k\times\mathsf H^k$ 求解 Hodge Laplace 问题 \eqref{eq:hodge-mixed-formulation}, $(\sigma_h,u_h,p_h)\in V_h^{k-1}\times V_h^k\times\mathsf H_h^k$ 求解离散问题 \eqref{eq:mixed-discrete-formulation}. 则}
\begin{equation}\label{eq:mixed-improved-dsigma}
\|d(\sigma-\sigma_h)\|\leq CE(d\sigma),
\end{equation}
\begin{equation}\label{eq:mixed-improved-sigma}
\|\sigma-\sigma_h\|\leq C[E(\sigma)+\eta E(d\sigma)],
\end{equation}
\begin{equation}\label{eq:mixed-improved-p}
\|p-p_h\|\leq C[E(p)+\mu E(d\sigma)],
\end{equation}
\begin{equation}\label{eq:mixed-improved-du}
\|d(u-u_h)\|\leq C(E(du)+\eta[E(d\sigma)+E(p)]),
\end{equation}
\begin{equation}\label{eq:mixed-improved-u}
\begin{aligned}
\|u-u_h\|&\leq C\big(E(u)+\eta[E(du)+E(\sigma)]\\
&\qquad+(\eta^2+\delta)[E(d\sigma)+E(p)]+\mu E(P_{\mathcal B}u)]\big),
\end{aligned}
\end{equation}
\BT{where the constant $C$ depends on the $W$-bounds on $\pi_h^{k-1}$, $\pi_h^k$, and $\pi_h^{k+1}$ but otherwise is independent of the choice of subcomplex.}{其中常数 $C$ 依赖于 $\pi_h^{k-1}$, $\pi_h^k$ 和 $\pi_h^{k+1}$ 的 $W$-界, 除此之外与子复形的选取无关.}
\end{Theorem}

% block: mixed-improved-error-proof-location page:72 type:ordinary-paragraph
\BT{The proof of this theorem is given in \BookCite{13}. It requires several pages and will not be repeated here.}{该定理的证明见 \BookCite{13}. 证明需要数页篇幅, 这里不再重复.}

% block: approximation-hilbert-summary-heading page:72 type:ordinary-paragraph
\Needspace{18\baselineskip}
\begin{tcolorbox}[enhanced,breakable,colback=KuangNei,frame hidden,boxrule=0pt,sharpish corners,before skip=10pt,after skip=10pt]
\refstepcounter{booksummary}\label{box:approximation-hilbert-summary}
\noindent\textbf{Box~\thebooksummary:} \BT{Summary of Chapter~\ref{chap:approximation-hilbert-complexes} on approximation of Hilbert complexes.}{Chapter~\ref{chap:approximation-hilbert-complexes} Hilbert 复形的逼近内容摘要.}\par

% block: approximation-hilbert-summary-discretization page:72 type:ordinary-paragraph
\BT{If $V^{k-1}\xrightarrow{d}V^k\xrightarrow{d}V^{k+1}$ is a segment of the domain complex of a closed Hilbert complex, we approximate the associated Hodge Laplacian problem by discretizing the mixed weak formulation via Galerkin's method with finite dimensional subspaces $V_h^j\subset V^j$.}{若 $V^{k-1}\xrightarrow{d}V^k\xrightarrow{d}V^{k+1}$ 是闭 Hilbert 复形的定义域复形的一个片段, 则使用有限维子空间 $V_h^j\subset V^j$, 通过 Galerkin 方法离散化混合弱表述, 来逼近相应的 Hodge Laplacian 问题.}

% block: approximation-hilbert-summary-requirements page:72 type:ordinary-paragraph
\BT{The key requirements on the Galerkin subspaces are (1) they afford good approximation, (2) they form a subcomplex ($dV_h^j\subset V_h^{j+1}$), and (3) they admit a \emph{bounded cochain projection}, i.e., bounded linear projections $V^j\to V_h^j$ which commute with the $d^j$.}{对 Galerkin 子空间的关键要求是 (1) 它们提供良好的逼近, (2) 它们构成子复形 ($dV_h^j\subset V_h^{j+1}$), (3) 它们具有\emph{有界上链投影}, 即与 $d^j$ 交换的有界线性投影 $V^j\to V_h^j$.}

% block: approximation-hilbert-summary-consequences page:72 type:ordinary-paragraph
\BT{If these requirements are satisfied, then the discrete complex is itself a Hilbert complex with the same cohomology, the space of discrete harmonic forms is close to that of the true harmonic forms, and the Poincaré constant for the discrete complex is bounded by the true Poincaré constant. It follows that the Galerkin method is consistent and stable and converges with the rate of the best approximation. Assuming also a compactness condition and boundedness of the cochain projections in $W$ norm, we further obtain optimal order error estimates for each variable separately with sharp regularity requirements.}{若这些要求得到满足, 则离散复形本身就是具有相同上同调的 Hilbert 复形, 离散调和形式空间接近真正的调和形式空间, 并且离散复形的 Poincaré 常数由真正的 Poincaré 常数控制. 由此可知, Galerkin 方法相容且稳定, 并按最佳逼近的阶收敛. 若再假设紧性条件以及上链投影在 $W$ 范数下有界, 则还能在精确的正则性要求下, 分别对各变量得到最优阶误差估计.}
\end{tcolorbox}

